\chapter{The Jacobian as an abelian variety}
\label{chap:Jacobian}

This short chapter isolates the \emph{Albanese universal property} of the identity
component \(J := \Pic^0_{C/k}\) of the Picard scheme of a smooth proper geometrically
irreducible curve \(C/k\) of genus \(g \ge 1\): for a fixed \(k\)-rational point
\(P \in C(k)\), the Abel--Jacobi morphism \(f^P\colon C \to J\) exhibits \(J\) as the
Albanese object of the pointed curve \((C,P)\).  This is Milne, \emph{Abelian Varieties},
III~\S6, Proposition~6.1.  Its proof is the load-bearing consumer of the
rational-map-extension Theorem~3.2 (\cref{thm:rational_map_to_av_extends}) and the
Milne~\S I.3 rigidity chain developed in \cref{chap:AbelianVarietyRigidity}; the
construction of \(\Pic^0_{C/k}\) itself, the symmetric-power / Abel--Jacobi machinery,
and the genus-stratified witness packaging live outside this subproject and are
referenced here only as context.

\section{The Albanese universal property of \(\Pic^0_{C/k}\)}

\begin{theorem}
[Albanese universal property of \(\Pic^0_{C/k}\) (Milne III Prop 6.1)]
  \label{thm:albanese_universal_property}
  
  \lean{AlgebraicGeometry.Scheme.Pic.albaneseUP}
  % SOURCE: [Milne, Abelian Varieties], Proposition 6.1, \S III.6, doc p.~104 (PDF p.~110) (read from references/abelian-varieties.pdf, PDF page 110)
  % SOURCE QUOTE: "PROPOSITION 6.1. Let \(P\) be a \(k\)-rational point on \(C\).
  % The map \(f^P\colon C\to J\) has the following universal property: for any
  % map \(\varphi\colon C\to A\) from \(C\) into an abelian variety sending \(P\)
  % to \(0\), there is a unique homomorphism \(\psi\colon J\to A\) such that
  % \(\varphi=\psi\circ f^P\)."
  % SOURCE QUOTE (Prop 6.4, the symmetric variant invoked through Cor 1.5;
  % read from references/abelian-varieties.pdf, PDF page 111): "PROPOSITION
  % 6.4. Let \(A\) be an abelian variety over \(k\). For any map
  % \(\varphi\colon C\times C\to A\) such that \(\varphi(\Delta)=0\), there is a
  % unique homomorphism \(\psi\colon J\to A\) such that \(\psi\circ F=\varphi\)."
  % SOURCE QUOTE (Remark 6.5, the Albanese characterization; read from
  % references/abelian-varieties.pdf, PDF page 111): "REMARK 6.5. The
  % proposition says that \((A,F)\) is the Albanese variety of \(C\) in the
  % sense of Lang 1959, II 3, p45. Clearly the pairs \((J,f^P)\) and \((J,F)\)
  % are characterized by the universal properties in (6.1) and (6.4)."
  \textit{Source: Milne, Abelian Varieties, III \S 6, Proposition 6.1 (p.~104), with Proposition 6.4 and Remark 6.5 (p.~105).}
  Let \(C\) be a smooth proper geometrically irreducible curve of genus \(g \geq 1\) over \(k\), fix a \(k\)-rational point \(P \in C(k)\), write \(J := \Pic^0_{C/k}\) for the identity component of its Picard scheme, and let \(f^P \colon C \to J\), \(Q \mapsto [\mathcal O_C(Q - P)]\), be the Abel--Jacobi morphism. For any morphism \(\varphi \colon C \to A\) from \(C\) to an abelian variety \(A\) over \(k\) with \(\varphi(P) = \eta_A\), there exists a unique homomorphism \(\psi \colon J \to A\) of group schemes such that \(\varphi = \psi \circ f^P\). Equivalently, \(J\) together with \(f^P\) exhibits \(J\) as the Albanese object of the pointed curve \((C, P)\) in the sense of \texttt{def:IsAlbanese}.
\end{theorem}
% SOURCE QUOTE PROOF: "PROOF. Consider the map
% \((P_1,\dots,P_g)\mapsto\sum_i\psi(P_i)\colon C^g\to A\). Clearly this is
% symmetric, and so it factors through \(C^{(g)}\). It therefore defines a
% rational map \(\psi\colon J\to A\), which (I 3.2) shows to be a regular map.
% It is clear from the construction that \(\psi\circ f^P=\varphi\) (note that
% \(f^P\) is the composite of \(Q\mapsto Q+(g-1)P\colon C\to C^{(g)}\) with
% \(f^{(g)}\colon C^{(g)}\to J\)). In particular, \(\psi\) maps \(0\) to \(0\), and
% (I 1.2) shows that it is therefore a homomorphism. If \(\psi'\) is a second
% homomorphism such that \(\psi'\circ f^P=\varphi\), then \(\psi\) and \(\psi'\)
% agree on \(f^P(C)+\dots+f^P(C)\) (\(g\) copies), which is the whole of \(J\)."
% [Milne, Prop 6.1, \S III.6, doc p.~104.]
\begin{proof}
  \textit{Source: Milne, Abelian Varieties, III \S 6, Proposition 6.1.}
  Symmetrise the \(g\)-fold sum: the morphism \((Q_1, \dots, Q_g) \mapsto \sum_i \varphi(Q_i) \colon C^g \to A\) is symmetric in the factors and so factors through the \(g\)-fold symmetric product \(C^{(g)} := C^g / S_g\). Composing with the inverse of the Abel--Jacobi map \(f^{(g)} \colon C^{(g)} \to J\), which is birational by Brill--Noether--Riemann--Roch (its generic fibre is a single point), one obtains a \emph{rational} map \(\psi \colon J \dashrightarrow A\). By Milne Theorem~3.2 (\cref{thm:rational_map_to_av_extends}: a rational map from a nonsingular variety to an abelian variety is everywhere defined --- the composite of Theorem~3.1's everywhere-extension-outside-codimension-\(2\) via the valuative criterion with Lemma~3.3's pure-codimension-\(1\) indeterminacy for maps into a group variety) the rational map \(\psi\) extends uniquely to a regular morphism \(J \to A\). The identity \(\psi \circ f^P = \varphi\) holds by construction once \(f^P\) is identified as the composite \(Q \mapsto Q + (g - 1)P \colon C \to C^{(g)}\) followed by \(f^{(g)} \colon C^{(g)} \to J\). The pointed condition \(\psi(0_J) = \eta_A\) then promotes \(\psi\) to a group-scheme homomorphism by Milne Corollary~1.2 (\texttt{lem:av_regular_map_is_hom}: a regular map between abelian varieties sending the identity to the identity is automatically a homomorphism). Uniqueness: any two such homomorphisms agree on \(f^P(C)\) and therefore on the sumset \(f^P(C) + \dots + f^P(C)\) (\(g\) summands), which fills \(J\) because \(f^{(g)} \colon C^{(g)} \to J\) is surjective (this is the dimension count \(\dim C^{(g)} = g = \dim J\) together with properness of \(C^{(g)}\)). The companion statement Proposition~6.4 for \(\varphi \colon C \times C \to A\) vanishing on the diagonal \(\Delta\) reduces to Proposition~6.1 via the additive decomposition Corollary~1.5 (\texttt{lem:hom_additivity_over_product}): the additivity decomposition \(\varphi = \varphi_1 \circ p_1 + \varphi_2 \circ p_2\) with \(\varphi_i \colon C \to A\), \(\varphi_i(P) = 0\), combined with \(\varphi|_\Delta = 0\), forces \(\varphi_1 = -\varphi_2\), after which Proposition~6.1 supplies the unique \(\psi\).
  \medskip
  \emph{Dependency audit (iter-172, Outcome (b)).} The chain of dependencies of this proof, classified by their formalisation status in the project's tree, is:
  \begin{itemize}
    \item \emph{In-tree, axiom-clean (no new work):} Corollary~1.5 (\texttt{lem:hom_additivity_over_product}, \texttt{AlgebraicGeometry.hom\_additive\_decomp\_of\_rigidity}); Corollary~1.2 (\texttt{lem:av_regular_map_is_hom}, \texttt{AlgebraicGeometry.av\_regularMap\_isHom\_of\_zero}); the Rigidity Lemma (\cref{thm:rigidity_lemma}) underlying both.
    \item \emph{Reserved Lean target, body not yet closed --- Mathlib gap on the codimension-\(1\) step:} Milne Theorem~3.2 (\cref{thm:rational_map_to_av_extends}, \texttt{AlgebraicGeometry.rationalMap\_to\_av\_extends}). The everywhere-extension-outside-codimension-\(2\) half (Theorem~3.1) is within Mathlib's reach via \texttt{ValuativeCriterion}; the pure-codimension-\(1\) half (Lemma~3.3, Weil-divisor / Auslander--Buchsbaum-flavoured argument that the indeterminacy locus of a map into a group variety is pure codimension~\(1\)) has \emph{no} Mathlib Weil-divisor API and is the route's riskiest sub-build (see \cref{rmk:thm32_codim1_mathlib_gap}). An open question for the prover: whether a \emph{pointwise} valuative extension --- run at each codimension-\(1\) point's discrete valuation ring separately --- side-steps building Weil-divisor theory entirely.
    \item \emph{Mathlib gap, shared with Route B:} symmetric powers of schemes \(C^{(g)} := C^g / S_g\) (Mathlib has no usable scheme-level quotient by a finite group action with the smooth-target property; the historical Route B was blocked on the same item). The Abel--Jacobi surjection \(f^{(g)} \colon C^{(g)} \to J\) and the birationality \(f^{(g)}\) uses Brill--Noether--Riemann--Roch, which routes through the cohomology layer common to Route~A sub-steps (A.1)--(A.3).
    \item \emph{Route A.3 output, consumed here:} the identity component \(J = \Pic^0_{C/k}\) together with its smoothness of dimension \(g\) and the surjectivity / dimension count for \(f^{(g)}\).
    \item \emph{NOT a dependency of A.4 (despite a textbook detour suggesting it might be):} autoduality \(J \cong J^\vee\) and the theorem of the cube. Milne's proof of Prop 6.1 as written does NOT invoke autoduality (which is established as Theorem 6.6 of the same section, \emph{after} 6.1, via the theorem of the cube). A pure Picard-functor / Poincar\'e-sheaf detour around Theorem~3.2 would route through autoduality and the cube, both of which the project has deliberately excised from the genus-\(0\) path (iter-163; \cref{rmk:cube_not_needed}). A.4 therefore commits to the Theorem~3.2 sub-build rather than the cube.
  \end{itemize}
  Per \cref{rmk:thm32_codim1_mathlib_gap}, the pure-codimension-\(1\) step (Lemma~3.3) requires either a Mathlib Weil-divisor / Auslander--Buchsbaum sub-build (a stand-alone Mathlib-prerequisite cascade in its own right, not currently scheduled) or the pointwise-valuative detour mentioned above. The STRATEGY.md A.4 row claim ``no new Mathlib namespace'' is incorrect with respect to this Lemma~3.3 sub-build; the planner should re-estimate A.4 with the codim-\(1\) extension treated as an explicit upstream piece.
\end{proof}
